% ============================================================
% Section 29: 一维谐振子与位移算符
% ============================================================
\section{量子力学：一维谐振子与位移算符——从对易关系到相干态}

\begin{modelbox}{题目}
一维谐振子，质量 $m$，频率 $\omega$。位置算符
\[
  X = \sqrt{\frac{\hbar}{2m\omega}}\,(a + a^\dagger).
\]
\textbf{(1)} 证明对任意解析函数 $F(X)$：
\[
  [a, F(X)] = \sqrt{\frac{\hbar}{2m\omega}}\,F'(X), \qquad
  [a^\dagger, F(X)] = -\sqrt{\frac{\hbar}{2m\omega}}\,F'(X).
\]
\textbf{(2)} 取 $F(X)=e^{ikX}$，证明递推关系：
\[
  \langle n|e^{ikX}|0\rangle = \frac{\alpha}{\sqrt{n}}\,\langle n-1|e^{ikX}|0\rangle,
  \quad \alpha = ik\sqrt{\frac{\hbar}{2m\omega}}.
\]
\textbf{(3)} 计算 $\langle n|e^{ikX}|0\rangle$ 的显式表达式。
\textbf{(4)} 初始态 $|\psi(0)\rangle = e^{ikX}|0\rangle$，求 $\langle X\rangle$ 和 $\langle P\rangle$。
\textbf{(5)} 求 $t$ 时刻能量测量结果为 $E_n$ 的概率 $P_n$。
\end{modelbox}

\vspace{0.3em}
\begin{insightbox}{结论预告}
$e^{ikX}|0\rangle$ 是\textbf{相干态}（纯虚数参数 $\alpha=ik\sqrt{\hbar/2m\omega}$），对应动量空间的平移。能量概率服从\textbf{泊松分布}，平均量子数 $\bar{n}=\hbar k^2/(2m\omega)$。
\end{insightbox}


\subsection{对易关系 $[a, F(X)]$ 的证明}

\subsubsection{基本对易子}

\begin{derivationbox}{从 $[a,X]$ 出发}
利用 $[a, a^\dagger]=1$ 和 $[a,a]=0$：
\begin{align*}
  [a, X] &= \sqrt{\frac{\hbar}{2m\omega}}\,[a, a+a^\dagger]
         = \sqrt{\frac{\hbar}{2m\omega}}, \\
  [a^\dagger, X] &= \sqrt{\frac{\hbar}{2m\omega}}\,[a^\dagger, a+a^\dagger]
                 = -\sqrt{\frac{\hbar}{2m\omega}}.
\end{align*}
注意 $[a,X]$ 和 $[a^\dagger,X]$ 都是\textbf{c 数}（与 $X$ 对易）。
\end{derivationbox}

\subsubsection{幂级数展开}

对 $F(X)=\sum_k c_k X^k$，利用对易子的 Leibniz 规则：
\[
  [a, X^n] = [a,X]X^{n-1} + X[a,X]X^{n-2} + \cdots + X^{n-1}[a,X].
\]
由于 $[a,X]$ 是 c 数，可与任意 $X$ 交换：
\[
  [a, X^n] = n[a,X]X^{n-1} = \sqrt{\frac{\hbar}{2m\omega}}\,nX^{n-1} = \sqrt{\frac{\hbar}{2m\omega}}\,\frac{d}{dX}X^n.
\]

由线性性推广到任意 $F(X)$：
\begin{formulabox}{核心对易关系}
\[
  \boxed{[a, F(X)] = \sqrt{\frac{\hbar}{2m\omega}}\,F'(X)}, \qquad
  \boxed{[a^\dagger, F(X)] = -\sqrt{\frac{\hbar}{2m\omega}}\,F'(X)}.
\]
\end{formulabox}

\begin{tipbox}{考场速记}
$[a, X^n] = nX^{n-1}[a,X]$ 的成立条件是 $[a,X]$ 为 c 数。对一般算符 $A,B$，$[A, B^n] \neq nB^{n-1}[A,B]$，除非 $[A,B]$ 与 $B$ 对易。
\end{tipbox}


\subsection{递推关系的推导}

取 $F(X)=e^{ikX}$，由 (1) 得：
\[
  [a, e^{ikX}] = ik\sqrt{\frac{\hbar}{2m\omega}}\,e^{ikX} = \alpha\,e^{ikX},
  \quad \alpha \equiv ik\sqrt{\frac{\hbar}{2m\omega}}.
\]

考虑矩阵元 $\langle n|e^{ikX}|0\rangle$。利用 $a^\dagger|n-1\rangle = \sqrt{n}\,|n\rangle$，即 $\langle n| = \frac{1}{\sqrt{n}}\langle n-1|a$：
\begin{align*}
  \langle n|e^{ikX}|0\rangle
  &= \frac{1}{\sqrt{n}}\langle n-1|a e^{ikX}|0\rangle \\
  &= \frac{1}{\sqrt{n}}\langle n-1|\bigl(e^{ikX}a + [a,e^{ikX}]\bigr)|0\rangle \\
  &= \frac{1}{\sqrt{n}}\langle n-1|\alpha e^{ikX}|0\rangle
   = \frac{\alpha}{\sqrt{n}}\langle n-1|e^{ikX}|0\rangle.
\end{align*}
（其中 $a|0\rangle=0$ 消去了第一项。）

\begin{formulabox}{递推关系}
\[
  \boxed{\langle n|e^{ikX}|0\rangle = \frac{\alpha}{\sqrt{n}}\,\langle n-1|e^{ikX}|0\rangle}
\]
逐次递推得：
\[
  \langle n|e^{ikX}|0\rangle = \frac{\alpha^n}{\sqrt{n!}}\,\langle 0|e^{ikX}|0\rangle.
\]
\end{formulabox}


\subsection{计算具体矩阵元}

\subsubsection{基态矩阵元 $\langle 0|e^{ikX}|0\rangle$}

利用谐振子基态波函数：
\[
  \psi_0(x) = \Bigl(\frac{m\omega}{\pi\hbar}\Bigr)^{1/4} e^{-m\omega x^2/2\hbar},
\]

\begin{align*}
  \langle 0|e^{ikX}|0\rangle
  &= \int_{-\infty}^{\infty} |\psi_0(x)|^2 e^{ikx}\,dx \\
  &= \sqrt{\frac{m\omega}{\pi\hbar}} \int_{-\infty}^{\infty} e^{-m\omega x^2/\hbar} e^{ikx}\,dx.
\end{align*}

利用高斯积分公式 $\displaystyle\int_{-\infty}^{\infty} e^{-ax^2+bx}dx = \sqrt{\frac{\pi}{a}}e^{b^2/4a}$（$a=m\omega/\hbar$，$b=ik$）：
\[
  \langle 0|e^{ikX}|0\rangle
  = \sqrt{\frac{m\omega}{\pi\hbar}} \cdot \sqrt{\frac{\pi\hbar}{m\omega}}\,
    e^{-\hbar k^2/4m\omega}
  = e^{-|\alpha|^2/2}.
\]

\begin{formulabox}{完整矩阵元}
\[
  \boxed{\langle n|e^{ikX}|0\rangle = \frac{\alpha^n}{\sqrt{n!}}\,e^{-|\alpha|^2/2}},
  \qquad \alpha = ik\sqrt{\frac{\hbar}{2m\omega}}.
\]
\end{formulabox}

\begin{insightbox}{这是相干态！}
标准的位移算符为 $D(\alpha)=\exp(\alpha a^\dagger - \alpha^* a)$。当 $\alpha$ 为\textbf{纯虚数}时，$\alpha^*=-\alpha$，于是
\[
  D(\alpha) = \exp\bigl(\alpha(a+a^\dagger)\bigr) = e^{ikX},
\]
其中 $k = \alpha\sqrt{2m\omega/\hbar}/i$。因此 $e^{ikX}|0\rangle$ 就是参数为纯虚数的相干态 $|\alpha\rangle$，它对应在\textbf{动量空间}中的平移（而非位置空间）。
\end{insightbox}


\subsection{$t=0$ 时的期望值}

\subsubsection{归一化验证}

\begin{align*}
  \langle\psi(0)|\psi(0)\rangle
  &= \langle 0|e^{-ikX}e^{ikX}|0\rangle = \langle 0|0\rangle = 1.
\end{align*}

\subsubsection{位置与动量期望值}

利用 $X = \sqrt{\frac{\hbar}{2m\omega}}(a+a^\dagger)$ 和 $P = i\sqrt{\frac{m\omega\hbar}{2}}(a^\dagger-a)$，在相干态 $|\alpha\rangle$ 中：
\[
  \langle X\rangle = \sqrt{\frac{\hbar}{2m\omega}}(\alpha+\alpha^*), \qquad
  \langle P\rangle = i\sqrt{\frac{m\omega\hbar}{2}}(\alpha^*-\alpha).
\]

由于 $\alpha = ik\sqrt{\frac{\hbar}{2m\omega}}$ 为\textbf{纯虚数}，$\alpha^*=-\alpha$：
\begin{align*}
  \langle X\rangle &= \sqrt{\frac{\hbar}{2m\omega}}(\alpha-\alpha) = \boxed{0}, \\
  \langle P\rangle &= i\sqrt{\frac{m\omega\hbar}{2}}(-2\alpha)
   = -2i\alpha\sqrt{\frac{m\omega\hbar}{2}} \\
  &= -2i\Bigl(ik\sqrt{\frac{\hbar}{2m\omega}}\Bigr)\sqrt{\frac{m\omega\hbar}{2}}
   = \boxed{\hbar k}.
\end{align*}

\begin{insightbox}{物理图像}
$e^{ikX}$ 作为动量空间的平移算符，给粒子一个动量``踢''(kick)，使其获得动量 $\hbar k$。由于 $\alpha$ 纯虚，位置期望值不变，但动量期望值增加了 $\hbar k$。这与 $e^{ikX}$ 在动量表象中是乘法算符 $e^{ikx} \to p$ 平移 $\hbar k$ 的图像完全一致。
\end{insightbox}


\subsection{$t$ 时刻能量测量概率}

\subsubsection{初始展开}

\begin{align*}
  |\psi(0)\rangle = e^{ikX}|0\rangle
  &= \sum_{n=0}^{\infty} |n\rangle\langle n|e^{ikX}|0\rangle \\
  &= \sum_{n=0}^{\infty} c_n |n\rangle,
  \quad c_n = \frac{\alpha^n}{\sqrt{n!}}e^{-|\alpha|^2/2}.
\end{align*}

\subsubsection{时间演化与概率}

时间演化：$|\psi(t)\rangle = \sum_n c_n e^{-iE_n t/\hbar}|n\rangle$，$E_n=\hbar\omega(n+\frac12)$。

能量测量概率（与时间无关）：
\begin{align*}
  P_n &= |\langle n|\psi(t)\rangle|^2 = |c_n|^2 \\
  &= \frac{|\alpha|^{2n}}{n!}e^{-|\alpha|^2} \\
  &= \boxed{\frac{1}{n!}\Bigl(\frac{\hbar k^2}{2m\omega}\Bigr)^n
     \exp\Bigl(-\frac{\hbar k^2}{2m\omega}\Bigr)}.
\end{align*}

\begin{formulabox}{泊松分布}
\[
  \boxed{P_n = \frac{\bar{n}^n}{n!}e^{-\bar{n}}, \qquad
    \bar{n} = \frac{\hbar k^2}{2m\omega} = |\alpha|^2}
\]
\end{formulabox}

\begin{insightbox}{为什么出现泊松分布？}
相干态 $|\alpha\rangle = D(\alpha)|0\rangle$ 是粒子数表象下的\textbf{最小不确定态}，其粒子数分布天然服从泊松分布。这是相干态的标志性特征：
\begin{itemize}
  \item 平均量子数 $\bar{n} = |\alpha|^2 = \frac{\hbar k^2}{2m\omega}$
  \item 方差 $(\Delta n)^2 = \bar{n}$（泊松分布的涨落特征）
  \item 能量期望值 $\langle E\rangle = \hbar\omega(\bar{n}+\frac12) = \frac{\hbar^2 k^2}{2m} + \frac{\hbar\omega}{2}$
\end{itemize}
其中 $\frac{\hbar^2 k^2}{2m}$ 正是动量平移 $\hbar k$ 对应的经典动能。
\end{insightbox}


\subsection{常见错误与考场策略}

\begin{errorbox}{易错点}
\begin{enumerate}
  \item \textbf{对易子与 $X$ 的交换性}：$[a,X^n]=nX^{n-1}[a,X]$ 仅在 $[a,X]$ 为 c 数时成立。若误用于 $[X,P^n]$ 会得到错误结果，因为 $[X,P]=i\hbar$ 不是 c 数与 $P$ 对易。
  \item \textbf{$\alpha$ 的虚实判断}：$\alpha=ik\sqrt{\hbar/2m\omega}$ 是纯虚数，因此 $\alpha^*=-\alpha$。若误用 $\alpha^*=\alpha$（实数情况），会得到 $\langle P\rangle=0$ 的错误结论。
  \item \textbf{能量概率的时间依赖性}：虽然态随时间演化，但 $|c_n e^{-iE_n t/\hbar}|^2=|c_n|^2$，概率\textbf{不随时间变化}。这是能量本征态测量的基本性质。
  \item \textbf{高斯积分的系数}：$\int e^{-ax^2+bx}dx = \sqrt{\pi/a}\,e^{b^2/4a}$ 中的 $a$ 在指数中是正数；注意归一化系数 $\sqrt{m\omega/\pi\hbar}$ 与积分结果 $\sqrt{\pi\hbar/m\omega}$ 恰好抵消。
\end{enumerate}
\end{errorbox}

\begin{cheatbox}{位移算符与相干态：速查卡}
\textbf{标准位移算符}：$D(\alpha)=\exp(\alpha a^\dagger - \alpha^* a)$

\textbf{位置空间位移}：$\alpha$ 为实数，$D(\alpha)=\exp(-i\frac{p_0}{\hbar}X)$，$\langle X\rangle = \sqrt{\frac{2\hbar}{m\omega}}\alpha$，$\langle P\rangle=0$

\textbf{动量空间位移}：$\alpha$ 为纯虚数，$D(\alpha)=\exp(i\frac{x_0}{\hbar}P)=e^{ikX}$，$\langle X\rangle=0$，$\langle P\rangle=\hbar k$

\textbf{通用公式}：
\[
  \langle X\rangle = \sqrt{\frac{\hbar}{2m\omega}}(\alpha+\alpha^*), \quad
  \langle P\rangle = i\sqrt{\frac{m\omega\hbar}{2}}(\alpha^*-\alpha)
\]
\textbf{能量概率}：始终泊松分布 $P_n = \bar{n}^n e^{-\bar{n}}/n!$，$\bar{n}=|\alpha|^2$
\end{cheatbox}
